\section{Introduction}\label{sec:intro}

Training efficiency is a central concern in deep learning for computer vision.
As models grow larger and datasets more complex, reducing the wall-clock time to reach a target accuracy directly lowers research iteration cost and energy consumption.
This problem spans two important settings: \emph{classification}, which assigns an image to a learned class, and \emph{reconstruction}, which predicts masked or corrupted regions of an image.
Both settings have seen rapid architectural progress, from convolutional neural networks (CNNs) to Vision Transformers (ViTs)~\cite{dosovitskiy2021vit}, yet efficient training recipes for these newer architectures remain underdeveloped.

A key gap exists in the transfer of training acceleration techniques.
Garg~\etal~\cite{garg2024cifar10} proposed four methods that dramatically accelerate CNN training on \cifar{}---frozen patch-whitening initialisation, Dirac (identity) initialisation, a ScaleBias output layer with the Lookahead optimiser, and alternating horizontal flip---achieving $94\%$ accuracy in 3.29 seconds on a single {GPU}.
However, it is unknown whether these methods carry their benefits to \emph{transformer} architectures or to \emph{reconstruction} tasks, where the inductive biases and training dynamics differ substantially from CNN classification.
Closing this gap is practically important: practitioners routinely apply CNN-era training recipes to ViTs without systematic validation, potentially wasting or even harming training efficiency.

We address this gap through a controlled empirical study.
We ask: \emph{Do the four CNN speed-up tricks of Garg~\etal{} transfer to (a) ViT-based classification and (b) CNN-based masked reconstruction?}
Our main contributions are:
\begin{enumerate}[leftmargin=*,topsep=2pt,itemsep=1pt]
    \item Two ViT classification baselines on \cifar{}: a \textbf{weak recipe} ($61.7\%$, 100 epochs) and a \textbf{strong recipe} ($81.4\%$, 300 epochs, RandAugment + Mixup/CutMix + DropPath), establishing the accuracy floor against which each trick is measured.
    \item Per-trick ablations under both recipes, isolating the effect of each technique on training efficiency and final accuracy.
    \item An analysis of why the same trick can help under one recipe and hurt under another, identifying the \emph{frozen-whitening/augmentation conflict} as the key mechanism.
    \item A controlled CNN masked-reconstruction experiment that directly tests trick transferability to reconstruction tasks, confirming the conflict mechanism and showing that whitening's benefit is preserved when inputs are unaugmented.
\end{enumerate}

Our central finding is that trick transferability is a \emph{trick~$\times$~recipe} interaction, not a fixed property of the trick or the architecture.
Frozen patch whitening accelerates convergence under basic augmentation ($+1.3$\pp{}) but hurts under strong augmentation ($-4.4$\pp{}); soft Dirac init is the only trick net-positive at the strong recipe level ($+0.4$\pp{}).
Practitioners should evaluate each CNN trick under the full target training recipe and task before adopting it.

The paper is organised as follows. Section~\ref{sec:related} covers related work. Section~\ref{sec:data} describes the dataset and preprocessing. Section~\ref{sec:methods} details the model architecture and each adapted trick. Section~\ref{sec:experiments} presents results, ablations, and mechanistic analyses. Section~\ref{sec:conclusion} concludes.


